% ============================================================
%  MÓDULO 09A — Problema resolvido pela arquitetura e análise
%  comparada com sistemas e técnicas disponíveis no GitHub.
%  Parte VIII. Livre de laudas.
% ============================================================
\chapter{O Problema que o Core Resolve e a Arquitetura Comparada}

\roteirodidatico
{Comparar arquiteturas começa por formular o problema que se quer resolver. Sem pergunta comum e critérios explícitos, diferenças de vocabulário podem parecer diferenças de capacidade.}
{\emph{Dimensão} é um aspecto comparado; \emph{critério} é a regra aplicada; \emph{alternativa} é outro sistema ou método; \emph{limite} indica o que a comparação não cobre.}
{Leia a pergunta, examine cada mecanismo, confira as fontes primárias e então avalie forças e custos usando critérios equivalentes para todos os casos.}
{Não inferir superioridade geral a partir de um único recurso. A comparação só vale para dimensões, versões, evidências e condições declaradas.}
{As alternativas foram avaliadas na mesma tarefa e versão, com dados comparáveis, ou a conclusão deve ser tratada apenas como análise documental?}

\casoguiado{comparação sob uma pergunta comum}
{Para comparar duas ferramentas de busca, use a mesma pergunta, o mesmo conjunto de documentos e critérios definidos antes de observar as saídas: recuperação de fontes pertinentes, identificação correta e tempo de resposta. Repita casos suficientes para mostrar variação. Uma leitura de documentação pode comparar recursos declarados, mas não substitui esse ensaio nem autoriza um ranking geral.}

\section{Pergunta de partida}

\elemento{Qual o problema, afinal de contas?}{\metainfo{ID}{A-00} \hfill \metainfo{Método}{comparação por dimensão}}

\begin{conceito}
O Core responde a uma pergunta de desenho: \emph{como operar um ecossistema com
mais de duzentos agentes, dezenas de skills, sete servidores MCP e um pipeline
científico, de modo que toda decisão deixe rastro, toda alegação tenha limite e
tudo continue rodando na máquina do operador?} A arquitetura não tenta \emph{ser}
o melhor modelo do mundo; tenta tornar \emph{auditáveis} a coordenação, a memória,
a prova e a honestidade de centenas de agentes.
\end{conceito}

\begin{resultado}
Decomposição do problema em seis subproblemas que a arquitetura cobre — e as sedes:\par
\smallskip
{\footnotesize
\begin{tabularx}{\textwidth}{YY}
\toprule
Subproblema & Mecanismo-resposta \\
\midrule
coordenação de muitos agentes & orquestrador + Blackboard A2A (Parte II) \\
memória persistente do que já foi feito & MetaBus + registro de evolução \\
prova de que ``ficou pronto'' & tríade spec $\to$ código $\to$ teste (SDD/TDD) \\
honestidade das alegações & anti-overclaim + honest reviewer \\
extensibilidade sem caos & catálogo regenerado + skills + MCP + hooks \\
operação local/offline & modelo local prioritário + instalador (Parte IX) \\
\bottomrule
\end{tabularx}\par}
\end{resultado}

\begin{niveis}
\textbf{Intuição:} uma empresa de agentes: sem memória ela esquece, sem contrato ela
mentiria sem querer, sem portaria ela não controla quem entra.\\
\textbf{Implementação:} ``ecossistema'' aqui é \emph{configuração derivada + contrato
executável + avaliação honesta} — não um framework único.\\
\textbf{Formalização:} arquitetura \emph{anti-overclaim por construção}: separa
capacidade-observada de capacidade-sancionada e faz cada etapa ser verificável.
\end{niveis}

\begin{descricao}
Para cada elemento, a ficha compara quatro coisas: \emph{problema que resolve},
\emph{forças} (técnicas e práticas), \emph{fraquezas} (técnicas e práticas) e
\emph{alternativas} identificáveis no GitHub. As alternativas foram escolhidas
entre projetos conhecidos e de código aberto, citadas pelo nome do repositório;
não há promessa de paridade medida — apenas contraste de desenho.
\end{descricao}

\clearpage
\section{Elemento 1: Coordenação por Blackboard A2A}

\elemento{Coordenação por blackboard A2A + orquestrador}

\begin{funcao}
\textbf{Problema resolvido (minucioso):} com 215 subagentes, uma delegação
``direta por nome'' é insustentável: o orquestrador não pode conhecer a fundo cada
capacidade o tempo todo. O blackboard inverte a responsabilidade — o orquestrador
\emph{anuncia} (publicar tarefa com capacidades requeridas) e os agentes
\emph{se voluntariam} por capacidade declarada (Agent Card). Isso baixa o
acoplamento (emissor não conhece receptor) e transforma a escolha em um leilão
governado por confiança e carga.
\end{funcao}

\begin{resultado}
Forças e fraquezas:\par
\smallskip
{\footnotesize
\begin{tabularx}{\textwidth}{YY}
\toprule
Positivas (técnicas e práticas) & Negativas (técnicas e práticas) \\
\midrule
desacopla orquestrador e agentes; adiciona agente sem recompilar o núcleo;
voluntariado por capacidade real & o ``melhor'' agente só é escolhido se seu
Agent Card for bem descrito; sem modelo de prioridade formal entre voluntários \\
rastro A2A audita quem se candidatou e quem venceu; tolera ausências com
fail-closed & overhead de coordenação por evento; leilão não negocia custo nem
prazo (não há contratos de SLA) \\
\bottomrule
\end{tabularx}\par}
\end{resultado}

\begin{metodo}
\textbf{Alternativas no GitHub:} microsoft/AutoGen, langchain-ai/LangGraph
(orquestração por grafos), crewAIInc/crewAI (equipes fixas),
openai/openai-agents-python (triagem declarativa) e \pth{microsoft/agent_connect}
(protocolo A2A). \textbf{Fortaleza relativa:} simplicidade determinística.
\textbf{Fraqueza relativa:} um LangGraph oferece retomada de fluxo por
checkpoint e gateways detalhados — o blackboard aqui é mais leve, porém menos
expressivo para fluxos com DAG complexa.
\end{metodo}

\clearpage
\section{Elemento 2: Memória (MetaBus + Blackboard)}

\elemento{Memória em dois canais com superfície MCP}

\begin{funcao}
\textbf{Problema resolvido (minucioso):} agentes sem estado repetem trabalho. O
Core separa \emph{memória de trabalho} (Blackboard: tarefas abertas) de
\emph{memória metacognitiva} (MetaBus: o que já se fez e se aprendeu), evitando a
confusão típica de bibliotecas que misturam estado de tarefa com histórico.
A leitura por \pth{mci_get_memory(topic, limit)} dá contexto ao próximo ciclo e o
fluxo publicar-assinar propaga aprendizado sem acoplamento.
\end{funcao}

\begin{resultado}
Comparativo:\par
\smallskip
{\footnotesize
\begin{tabularx}{\textwidth}{YY}
\toprule
Positivas & Negativas \\
\midrule
dois canais com papéis distintos reduzem ruído; tópicos permitem assinatura
seletiva; exposição MCP deixa consumir por qualquer cliente & memória local a uma
máquina; sem versionamento de estado entre hosts \\
registro de reflexão (metacognition.reflected) alimenta decisões futuras &
sem políticas de expiração/compressão de tópicos automáticas em larga escala \\
\bottomrule
\end{tabularx}\par}
\end{resultado}

\begin{metodo}
\textbf{Alternativas no GitHub:} cpacker/MemGPT, letta/letta (memória
hierárquica), getzep/zep (memória temporal para agentes) e mem0ai/mem0.
\textbf{Fortaleza relativa:} desenho simples de dois canais e rastro A2A.
\textbf{Fraqueza relativa:} mem0 e letta oferecem \emph{recuperação semântica
escalável} com vetores dedicados; aqui a escala é a de um repositório local.
\end{metodo}

\clearpage
\section{Elemento 3: Prova por SDD/TDD}

\elemento{Especificação executável, evidência tipada}

\begin{funcao}
\textbf{Problema resolvido (minucioso):} em ecossistemas com muitos agentes, o
risco-mor é o \emph{verde falso} — uma entrega que ``parece pronta'' sem prova.
O SDD/TDD amarra spec a teste e exige dois artefatos de evidência
(\pth{TrustedTestEvidence} e \pth{CriterionRuntimeRecord}), com caminhos
canonizados e seletores seguros. O \textbf{runner} repete o teste sem aceitar
contagens ambíguas (fail-closed).
\end{funcao}

\begin{resultado}
Comparativo:\par
\smallskip
{\footnotesize
\begin{tabularx}{\textwidth}{YY}
\toprule
Positivas & Negativas \\
\midrule
rastreabilidade spec $\to$ teste auditável; evidência tipada (hashes, registros)
impede ``verde por intimidação''; gates determinísticos & a qualidade do critério
depende de quem redige a spec; sem orquestrador de suíte distribuída \\
modelagem de critérios executáveis e recusa a testes ambíguos (boas práticas de
engenharia de testes) & herda limitações do pytest/harness para flake e timing \\
\bottomrule
\end{tabularx}\par}
\end{resultado}

\begin{metodo}
\textbf{Alternativas no GitHub:} behave/cucumber e pytest-bdd (BDD/Gherkin),
hypothesis (testes por propriedade), confident-ai/DeepEval e
explodinggradients/ragas (avaliação de RAG). \textbf{Fortaleza relativa:} o
casamento da spec com as duas formas de evidência e a canonização de caminho.
\textbf{Fraqueza relativa:} BDD expressa comportamento de negócio e ecossistemas
como DeepEval têm métricas de qualidade de LLM que aqui precisam ser cobertas com
engines próprias, com menor cobertura pronta.
\end{metodo}

\clearpage
\section{Elemento 4: Escrita científica (MASWOS)}

\elemento{Pipeline argumentativo contra o "artigo-narrativa"}

\begin{funcao}
\textbf{Problema resolvido (minucioso):} redigir ciência em esteira de
especialistas (diagnóstico do escopo, busca, evidências, metodologia,
estatística, resultados, discussão, auditoria ABNT) impede que o texto seja uma
``história coerente'' sem estrutura verificável. A função argumentativa de cada
seção é projetada antes; o QA devolve o texto para revisão quando a função não se
sustenta.
\end{funcao}

\begin{resultado}
Comparativo:\par
\smallskip
{\footnotesize
\begin{tabularx}{\textwidth}{YY}
\toprule
Positivas & Negativas \\
\midrule
divisão explícita de trabalho argumentativo; auditoria de citações ABNT acoplada;
rastro MASWOS--agente por escala & não substitui revisão por pares nem leitor
humano de plágio; mais leve em geração autônoma de texto longo \\
boas práticas para TCC/tese/artigo com checklists auditáveis &
foco acadêmico brasileiro (ABNT) é vantagem local, menos útil para periódicos
em inglês sem adaptação \\
\bottomrule
\end{tabularx}\par}
\end{resultado}

\begin{metodo}
\textbf{Alternativas no GitHub:} stanford-oval/storm (geração de artigos com
pesquisa automática), \pth{binary-husky/gpt_academic} (assistência acadêmica) e
timpara/opencode-academic-research (skill acadêmico). \textbf{Fortaleza
relativa:} auditoria normativa e função argumentativa por estágio.
\textbf{Fraqueza relativa:} STORM automatiza síntese de múltiplas fontes em
escala; MASWOS depende de degraus orquestrados manualmente.
\end{metodo}

\clearpage
\section{Elemento 5: Auditores (scanners) e integridade}

\elemento{Múltiplos raios de auditoria com verificador de Merkle}

\begin{funcao}
\textbf{Problema resolvido (minucioso):} saber \emph{se} um texto (literário ou
científico) tem problemas é uma leitura multidimensional (narrativa, estilo,
símbolos, ética, teoria, rigor). Os scanners executam essas leituras em paralelo
e consolidam em escore; a integridade do próprio código-fonte é verificada por
\emph{árvore de Merkle} — o auditor audita, e também é auditável.
\end{funcao}

\begin{resultado}
Comparativo:\par
\smallskip
{\footnotesize
\begin{tabularx}{\textwidth}{YY}
\toprule
Positivas & Negativas \\
\midrule
auditoria multi-critério extensível; verificação de integridade por Merkle do
código & escores SRI/EXS são heurísticos, sem chancela externa; não detectam
plágio como um iThenticate \\
anti-overclaim embutido nos relatórios (o escore não se auto-absolutiza) &
cobertura de scanners próprios é a documentada, sem leque pronto de empacotadores \\
\bottomrule
\end{tabularx}\par}
\end{resultado}

\begin{metodo}
\textbf{Alternativas no GitHub:} astral-sh/ruff e ESLint (estilo em código; não
literário), Arize-ai/phoenix e UpTrain (observabilidade de LLM).
\textbf{Fortaleza relativa:} scanners literários e científicos num único
pipeline. \textbf{Fraqueza relativa:} Phoenix/LangFuse dão traços e telemetria
de sessão cross-host; aqui não há telemetria entre máquinas.
\end{metodo}

\clearpage
\section{Elemento 6: Avaliação honesta}

\elemento{Nota que se recusa a mentir}

\begin{funcao}
\textbf{Problema resolvido (minucioso):} avaliação de IA costuma inflar: um
``juiz-LLM'' premia coerência narrativa com mérito. O \pth{honest_reviewer}
separa \emph{cobertura} de \emph{mérito}, detecta \emph{inflação de alegações}
palavra a palavra e devolve uma \emph{faixa} em vez de nota única; sem validação
externa, a classe \pth{verified} permanece bloqueada.
\end{funcao}

\begin{resultado}
Comparativo:\par
\smallskip
{\footnotesize
\begin{tabularx}{\textwidth}{YY}
\toprule
Positivas & Negativas \\
\midrule
detecção determinística de inflação (termos-chave exigem chancela) &
faixa reflete cobertura local; não produz ranking comparativo entre sistemas \\
devolve banda de escore e limitações explícitas; guarda anti-overclaim &
molda-se pela política do Core; parâmetros de banda fixos \\
\bottomrule
\end{tabularx}\par}
\end{resultado}

\begin{metodo}
\textbf{Alternativas no GitHub:} openai/evals (juiz-LLM), confident-ai/DeepEval
(LLM-as-judge) e Arize-ai/phoenix evals. \textbf{Fortaleza relativa:} a rigidez
anti-inflação é rara e pragmática. \textbf{Fraqueza relativa:} evals, DeepEval
e Phoenix têm dezenas de métricas prontas (fidelidade, alucinação) e
infraestrutura de comparação que aqui não existem.
\end{metodo}

\clearpage
\section{Elemento 7: Confiança por ação}

\elemento{TrustScorer como filtro operacional}

\begin{funcao}
\textbf{Problema resolvido (minucioso):} sem histórico, todo agente reaparece
neutro. O \pth{TrustScorer} guarda confiança por \emph{ação}: cada resultado
atualiza o escore (\pth{record_outcome}), a deriva é avaliada e os gates
operacionais (confiança mínima e carga máxima) decidem quem pode assumir o quê.
\end{funcao}

\begin{resultado}
Comparativo:\par
\smallskip
{\footnotesize
\begin{tabularx}{\textwidth}{YY}
\toprule
Positivas & Negativas \\
\midrule
leve, determinístico, sem rede; por-ação mantém granularidade; deriva captura
mudança de comportamento & heurística sem garantia probabilística; sem meio de
exploração formal (bandit) — só atualiza escore \\
aplicável em runtime com custo desprezível & não agregado por domínio
automaticamente; parâmetro de limiar fixo \\
\bottomrule
\end{tabularx}\par}
\end{resultado}

\begin{metodo}
\textbf{Alternativas no GitHub:} langfuse/langfuse (observabilidade) e
bibliotecas de bandidos (MAB). \textbf{Fortaleza relativa:} custo mínimo e
decisões em runtime com determinismo. \textbf{Fraqueza relativa:} LangFuse dá
traços de custo/latência por chamada — ausentes aqui — e bandidos equilibram
exploração-explotação que este escore não faz.
\end{metodo}

\clearpage
\section{Elemento 8: Registro de evolução}

\elemento{Append-only com âncora de Merkle}

\begin{funcao}
\textbf{Problema resolvido (minucioso):} ``o que mudou e por quê'' fica na
burocracia humana. O \pth{EvolutionRegistry} grava cada mudança como ciclo com
\emph{hashes de artefatos} e \emph{encadeamento por Merkle}: alterar o passado
quebra a corrente. Ciclos \emph{recusados} também ficam registrados — a memória
do que não deu certo é tão valiosa quanto a de acertos.
\end{funcao}

\begin{resultado}
Comparativo:\par
\smallskip
{\footnotesize
\begin{tabularx}{\textwidth}{YY}
\toprule
Positivas & Negativas \\
\midrule
histórico à prova de adulteração local; registra rejeições (memória negativa) &
formato próprio, não interoperável com SLSA/in-toto/sigstore \\
numeração sequencial e ancoragem auditáveis; integrado ao ciclo de reflexão &
sem assinatura por autoridade externa; cadeia localmente mantida \\
\bottomrule
\end{tabularx}\par}
\end{resultado}

\begin{metodo}
\textbf{Alternativas no GitHub:} in-toto/in-toto e sigstore (cadeia de
suprimento assinada) e slsa-framework/slsa. \textbf{Fortaleza relativa:}
simplicidade e foco no ciclo de desenvolvimento. \textbf{Fraqueza relativa:}
in-toto/sigstore dão prova criptográfica de suprimento contra terceiros — aqui a
cadeia é local e auto-assinada.
\end{metodo}

\clearpage
\section{Elemento 9: Roteamento de modelos}

\elemento{Local primeiro, nuvem com fallback}

\begin{funcao}
\textbf{Problema resolvido (minucioso):} escolher modelo por \emph{tipo de
tarefa} e por \emph{soberania}: prioridade local (LiteRT/Colibri/Ollama/PAIR)
preserva privacidade e segue operando sem rede; quando falta, degrada para API
com \pth{fallback_used=True} — sempre transparente.
\end{funcao}

\begin{resultado}
Comparativo:\par
\smallskip
{\footnotesize
\begin{tabularx}{\textwidth}{YY}
\toprule
Positivas & Negativas \\
\midrule
privacidade por omissão (dados não saem da máquina); fallback explícito; perfis
de tarefa claros & sem balanceamento de custo por requisição; sem cache de
respostas \\
configuração por perfil (R607) documentada; PAIR opcional sem exceções
silenciosas & menos expressivo que gateways maduros para produção multiusuário \\
\bottomrule
\end{tabularx}\par}
\end{resultado}

\begin{metodo}
\textbf{Alternativas no GitHub:} BerriAI/litellm (unificação de providers),
openrouter (federação de modelos) e songquanpeng/one-api (gateway unificado).
\textbf{Fortaleza relativa:} local-first com opção de PAIR e medição local
documentada. \textbf{Fraqueza relativa:} LiteLLM e one-api suportam budget,
rate-limit e balanceamento por crédito — capacidades ausentes aqui.
\end{metodo}

\clearpage
\section{Elemento 10: Simulação social (MiroFish)}

\elemento{Opinião simulada determinística}

\begin{funcao}
\textbf{Problema resolvido (minucioso):} antever reação pública exige iterar
cenários. O MiroFish local é enxame determinístico (vieses cíclicos, validação
cruzada); o MiroFish-Offline (AGPL) entra por composição com fail-closed. Toda a
operação é sintética, o que permite explorar ``e se a comunicação for assim?''
sem o custo ético-operacional de testar em humanos reais.
\end{funcao}

\begin{resultado}
Comparativo:\par
\smallskip
{\footnotesize
\begin{tabularx}{\textwidth}{YY}
\toprule
Positivas & Negativas \\
\midrule
determinístico e de baixa dependência no motor local; integração fail-closed;
licença respeitada por composição & escala de agentes modesta frente a frameworks
dedicados; sem grafos de comportamento sofisticados \\
relatório estruturado para interpretação humana; rodapé ``não é medição de
opinião real'' & resultado nunca valida a realidade — serve só como cenário \\
\bottomrule
\end{tabularx}\par}
\end{resultado}

\begin{metodo}
\textbf{Alternativas no GitHub:} projectmesa/mesa (agent-based em larga escala),
alibaba/AgentScope e cameldeng/COALA - CAMEL-AI (orquestração multiagente) e
nikmcfly/MiroFish-Offline (fonte deste). \textbf{Fortaleza relativa:} espelho de
comunicação e baixa dependência. \textbf{Fraqueza relativa:} mesa e AgentScope
oferecem visualizações, ticks paralelos e milhões de execuções; aqui a herança é
pequena.
\end{metodo}

\clearpage
\section{Elemento 11: Integração por MCP}

\elemento{Protocolo em vez de plugins proprietários}

\begin{funcao}
\textbf{Problema resolvido (minucioso):} cada ferramenta com API própria vira
incompatível. O MCP padroniza ferramentas entre CLI, skills, plugins e agentes:
7 servidores expõem capacidade reutilizável, e servidores nativos
(\pth{.opencode/mcp/}) mostram o Core atuando como fornecedor de protocolo.
\end{funcao}

\begin{resultado}
Comparativo:\par
\smallskip
{\footnotesize
\begin{tabularx}{\textwidth}{YY}
\toprule
Positivas & Negativas \\
\midrule
padrão aberto entre autenticação e contrato de ferramenta; extensível com MCP-Go
e SDKs várias & custo de mensageria/tokens por chamada; dependência de um cliente
MCP bem implementado \\
integrantes externos entram sem conhecimento interno do repositório &
superfície protocolada mas sem catálogo de descoberta ativa (MCP-Zero) além do
configurado \\
\bottomrule
\end{tabularx}\par}
\end{resultado}

\begin{metodo}
\textbf{Alternativas no GitHub:} modelcontextprotocol (mcp, SDKs) versus
ecossistemas plugin-only (ex.: plugins TypeScript do OpenCode, que coexistem no
Core). \textbf{Fortaleza relativa:} padrão de comunicação com futuro aberto.
\textbf{Fraqueza relativa:} descoberta ativa de ferramentas (MCP-Zero) reduziria
o overhead de tokens; aqui os 7 MCPs são fixos em configuração.
\end{metodo}

\clearpage
\section{Síntese da comparação}

\begin{resultado}
Mapa resumido (leitura rápida):\par
\smallskip
{\footnotesize
\begin{tabularx}{\textwidth}{Yl}
\toprule
Elemento & Vantagem diferencial sobre alternativas do GitHub \\
\midrule
Blackboard A2A & leilão voluntário determinístico, leve \\
MetaBus + Blackboard & dois canais com superfície MCP \\
SDD/TDD & evidência tipada + canonização segura \\
MASWOS & função argumentativa por estágio + auditoria ABNT \\
Scanners & multi-critério + Merkle do código \\
Honest reviewer & recusa automática de inflação de alegação \\
TrustScorer & confiança por ação, determinística \\
Evolution/Merkle & memória de rejeições auditável \\
Roteamento local & privacidade por omissão + fallback explícito \\
MiroFish & síntese determinística com fail-closed \\
MCP & protocolo aberto sem acoplamento proprietário \\
\bottomrule
\end{tabularx}\par}
\end{resultado}

\begin{limite}
Comparações são de \emph{desenho}, não de benchmark medido: nenhum destes
contrastes implica que o Core seja ``melhor'' ou ``pior'' em desempenho — apenas
que resolve os mesmos subproblemas com escolhas diferentes, cujas forças e
fraquezas este capítulo documenta com transparência.
\end{limite}

\section{Resumo e exercícios}

\begin{resultado}
\textbf{O que você deve ter aprendido:} (1) o problema central é operar muitos
agentes com memória, prova, honestidade e soberania local; (2) cada elemento tem
uma ficha de problema $\to$ forças $\to$ fraquezas $\to$ alternativas; (3)
comparação é um contraste de \emph{desenho}, nunca um veredito de desempenho.
\textbf{Pergunta de fixação:} escolha um elemento deste capítulo e replique a
ficha para um caso da sua própria área (ex.: um editor de código comparado ao
IDE que você usa). Em seguida, verifique se as fichas deste capítulo se aplicam
ao seu exemplo.
\end{resultado}

% ============================================================

%  Gerado por gen_mod14_ativacao.py (R613) — seção de ativação e cálculo
%  para o módulo mod-09a-comparada.tex. Edições manuais são preservadas nas próximas
%  execuções (marcador impede reescrita).
% ============================================================

\section[Leitura guiada]{Leitura guiada — Arquitetura comparada: comparar sem ranquear}

\elemento{Leitura guiada — Arquitetura comparada}{\metainfo{ID}{N-09a} \hfill \metainfo{Ciclo}{R616} \hfill \metainfo{Estilo}{comparar não é ranquear}}

\begin{descricao}
\textbf{A pergunta.} O que se pode concluir ao comparar dois sistemas? Uma
comparação arquitetural pode descrever diferenças de componentes, interfaces e
fluxos. Já uma afirmação de desempenho exige que os sistemas sejam medidos em
condições comparáveis. Este capítulo examina \textbf{11 elementos} do Core,
explicita problemas, escolhas e alternativas e pergunta, em cada caso, se existe
uma comparação externa com medidas reproduzíveis. Descrever uma diferença não é
ranquear as soluções.
\end{descricao}

\begin{definicao}
\textbf{Grandezas e definições.}
\begin{itemize}[leftmargin=1.3em,itemsep=1pt]
  \item \textbf{Elemento} ($e$): mecanismo do Core comparado; $e = 11$.
  \item \textbf{Cobertura comparativa externa} ($b$): elementos com comparação externa medida, dividido por $e$.
  \item \textbf{Desenho}: escolha estrutural (protocolo, canal, gate).
  \item \textbf{Desempenho}: medida observada sob condição declarada.
  \item \textbf{Afirmação comparativa de desempenho}: conclusão que exige $b > 0$, métrica comum e condições de medição comparáveis. Sem isso, ainda é possível descrever diferenças de arquitetura, mas não declarar superioridade empírica.
\end{itemize}
\end{definicao}

\begin{metodo}
\textbf{Dedução passo a passo.} O que seria necessário para um veredito legítimo?
\begin{enumerate}[leftmargin=1.4em,itemsep=2pt]
  \item Mesmo problema: as duas soluções precisam atender ao mesmo requisito, não a subproblemas diferentes.
  \item Mesma métrica e mesma condição: mesma tarefa, mesmo volume, mesmo hardware.
  \item Fonte verificável: código, publicação ou repositório do outro lado.
  \item Só então $b$ deixa de ser zero e o veredito pode ser formulado.
\end{enumerate}
\end{metodo}

\begin{resultado}
\textbf{Exemplo resolvido (a conta honesta).} Nesta conferência, o manual não
localizou medição externa comparável documentada para nenhum dos $11$ elementos:
\[ b = \frac{0}{11} = 0\%. \]
Ou seja: $b = 0$ e, portanto, não há base neste inventário para uma afirmação de
superioridade empírica. Isso não torna a afirmação falsa; significa que ela não
está sustentada pela comparação documentada aqui. Uma medição comparável em um
elemento elevaria a fração para $1/11 \approx 0{,}091$, ou cerca de $9{,}1\%$.
\end{resultado}

\begin{resultado}
\textbf{Ordem de grandeza e verificação.} Uma cobertura comparativa de $0/11$
($b = 0\%$) não é ``quase suficiente'': é ausência total de evidência
comparativa. Mesmo que a base de 11 implementações fosse pequena para um
\emph{julgamento} relativo, ela seria suficiente para a afirmação fraca
``não há evidência neste repositório de que X seja superior'' --- que é
justamente o limite que o texto sustenta. \emph{Verifique você mesmo:} liste os 11
sistemas citados no capítulo e, para cada um, procure no repositório um artefato de
comparação; a ausência esperada em todos confirma o próprio exemplo.
\end{resultado}

\begin{aviso}
\textbf{Limite de validade.} Este exemplo não mede o Core; mede a \emph{lacuna de
medição}. Ele existe para que o leitor não confunda contraste de desenho com
vantagem demonstrada --- o mesmo princípio anti-overclaim aplicado à arquitetura.
\end{aviso}

\begin{resultado}
\textbf{Exercícios.} (1) Escolha um dos 11 elementos e liste o que precisaria para
que $b$ deixasse de ser zero naquele item. (2) Reescreva, em uma frase sem superlativo, a vantagem de desenho do Blackboard A2A.
\end{resultado}

% ATIVACAO-R613
\section[Ativação e cálculo]{Ativação e cálculo — Comparação e benchmarking}
\elemento{ativ-09a}{\metainfo{Ciclo}{R613} \hfill \metainfo{Módulo}{mod-09a-comparada.tex}}

\begin{processo}
GATILHO. ``compare'', ``benchmark'', ``segunda opinião''.
ROTA. avaliação multi-agente; gate: ponderação igual e ablação.
FUNÇÃO. comparar agentes/motores com critérios justos e reproducíveis.
\end{processo}

\begin{resultado}
CÁLCULO. média ponderada igual entre geradores; ablação isola cada componente; seeds fixas garantem determinismo.
\end{resultado}

\begin{descricao}
JUSTIFICATIVA TEÓRICA. comparar exige reportar o processo — só há mérito comparável se o relato do método é completo. O lastro teórico vem da citação que segue: recorte conferido em 29/09/2026, com a página de origem e tradução livre e fiel.
\end{descricao}

\begin{aviso}
LIMITE E ANTI-OVERCLAIM. comparação interna não substitui avaliação por pares externa. A verificação atesta a existência e a
localização da fonte (R110), não o mérito da afirmação.
\end{aviso}

\noindent\textit{Interpretação e cálculo ilustrativo:} O artigo mostra que a força de um resultado depende da probabilidade prévia, do poder estatístico e do limiar de falso positivo; não permite concluir que qualquer achado específico seja falso. No modelo sem viés, $PPV=((1-\beta)R)/(R-\beta R+\alpha)$. Supondo apenas para ilustração $R=0{,}1$ e $\alpha=0{,}05$, um poder de $0{,}2$ ($\beta=0{,}8$) dá $PPV=0{,}02/(0{,}1-0{,}08+0{,}05)\approx0{,}286$; com poder $0{,}8$ ($\beta=0{,}2$), $PPV=0{,}08/(0{,}1-0{,}02+0{,}05)\approx0{,}615$. São premissas didáticas, não estimativas do benchmark do livro. A consequência prática é separar triagem interna de confirmação independente \citref{IOANNIDIS, J. P. A. Why most published research findings are false. \emph{PLoS Medicine}, v.~2, n.~8, e124, 2005}{10.1371/journal.pmed.0020124}{p.~e124 (resumo, p.~e124). é a âncora do ceticismo estruturado do benchmarking: a maioria dos achados publicados é falsa; comparar sem controle de vieses de desenho e relato reproduz o mesmo problema.}{There is increasing concern that most current published research findings are false}{Há preocupação crescente de que a maioria dos resultados de pesquisa atualmente publicados seja falsa.}{IOANNIDIS, 2005, p.~e124}
